\section{Independent spectral orders and independent shifts}
\label{sec:polarization}

\subsection{A mixed determinant with no mixed zeta products}
Let $(s_\nu,a_\nu)$, $1\le\nu\le r$, be independent pairs with
$\Rea s_\nu>1$ and $\Rea a_\nu>0$. The finite-difference order is
still $d=2r-2$ for every weight. Define the averaged symmetrization
\begin{equation}\label{eq:sympartition}
 \Z_r^{\sym}(\bm s,\bm a)
 =\frac1{r!}\sum_{\sigma\in S_r}
   \sum_{0\le n_1<\cdots<n_r}\Delta(\bm n)^2
      \prod_{i=1}^r w_{n_i}(s_{\sigma(i)},a_{\sigma(i)}).
\end{equation}
When all pairs coincide, this is exactly $\Z_r(s,a)$, not $r!\Z_r(s,a)$.

For independent formal variables $t_1,\ldots,t_r$, write
\[
 E(\bm t)=\sum_{\nu=1}^r t_\nu E(s_\nu,a_\nu),\qquad
 C(\bm t)=\det E(\bm t)_{0\le i,j<r-1}.
\]
The latter is a homogeneous polynomial of degree $r-1$.

\begin{theorem}[Polarized single-zeta closure]\label{thm:polarized}
The series \eqref{eq:sympartition} is absolutely and locally normally
convergent in the stated domain, and
\begin{align}
 \Z_r^{\sym}(\bm s,\bm a)
 ={}&\frac1{r!}[t_1\cdots t_r]\det E(\bm t)\notag\\
 &+\frac1{r!}\sum_{\nu=1}^r\zeta(s_\nu,a_\nu)
       \left[\prod_{\kappa\ne\nu}t_\kappa\right]C(\bm t).
       \label{eq:polarized}
\end{align}
The coefficient expressions are finite and entire in all spectral
variables. Thus the joint continuation contains only a linear combination
of single Hurwitz zeta functions. It has no products of zeta functions
and no product poles at intersections of $s_\nu=1$.
\end{theorem}
\begin{proof}
Each term in the finite symmetrization satisfies the product majorant from
Lemma~\ref{lem:decay}, using $\min_\nu\Rea s_\nu>1$ on compact subsets.
Apply Cauchy--Binet to the combined weight
$\sum_\nu t_\nu w_n(s_\nu,a_\nu)$. Its moment matrix is
\[
 E(\bm t)+\left(\sum_\nu t_\nu\zeta(s_\nu,a_\nu)\right)
       \ee\ee^{\mathsf T}.
\]
Its determinant is
\[
 \det E(\bm t)+
 \left(\sum_\nu t_\nu\zeta(s_\nu,a_\nu)\right)C(\bm t).
\]
On the series side, the coefficient of $t_1\cdots t_r$ sums all $r!$
assignments of the labeled weights to the ordered particles. Divide by
$r!$. This proves the formula and its precise normalization. The
meromorphic conclusion follows because the displayed finite coefficients
are entire in the independent spectral variables.
\end{proof}

\begin{remark}[Why independent shifts do not increase the rank]
Each separate shift creates its own Hurwitz zeta coordinate, but every
coordinate enters the same bottom-right moment entry. The number of
special-function coordinates and the rank of the matrix subspace carrying
them are different quantities. This is the mechanism behind the
multiple-residue-class extension in Section~\ref{sec:classification}.
\end{remark}

For $r=2$ and a common shift $a$, set
$A_s=a^{-s}$ and $B_s=(a+1)^{-s}$. Equation~\eqref{eq:polarized} reads
\begin{align}
 \Z_2^{\sym}(s,t;a)
 ={}&\frac{A_s-B_s}{4}\zeta(t,a)
     +\frac{A_t-B_t}{4}\zeta(s,a)\notag\\
 &-\frac{A_sA_t}{2}
     +\frac{A_sB_t+A_tB_s}{8}.
 \label{eq:polarized-r2}
\end{align}
In particular,
\begin{equation}\label{eq:mixed23}
 \Z_2^{\sym}(2,3;1)
   =\frac7{32}\zeta(2)+\frac3{16}\zeta(3)-\frac{29}{64}.
\end{equation}
The unaveraged two-term symmetrization is twice this value.

\subsection{All mixed order derivatives}
For every integer $q\ge0$,
\begin{equation}\label{eq:order-weight}
 \partial_s^q w_n(s,a)
 =\frac{(-1)^q}{d!}\sum_{j=0}^d(-1)^j\binom dj
       \frac{\log^q(n+a+j)}{(n+a+j)^s}.
\end{equation}
Differentiating \eqref{eq:polarized} in any fixed collection of the
independent $s_\nu$ therefore evaluates symmetrized sums with independent
finite-difference logarithmic weights. Their right sides contain only
finite coefficients times single derivatives
$\zeta^{(j)}(s_\nu,a_\nu)$.

This is not a formal differentiation of integer-weight identities.
The independent-variable identity is first established on a complex
product domain with normal convergence. Subsequent integer specializations
of its derivatives are consequently justified. Differentiating in the
shifts as well produces mixed logarithmic and polygamma identities
without changing that analytic foundation.

\subsection{A useful distinction from general determinant sums}
A generic determinant of zeta-valued moments contains products of zeta
functions. Here their absence has an exact reason: the coefficient
matrices of all the zeta coordinates share a one-dimensional image.
No arithmetic cancellation conjecture is being used. Conversely, the
formula does not say that an arbitrary unsymmetrized assignment of distinct
orders has the same reduction. The symmetrization is essential to this
particular determinant argument.
